\documentclass[11pt]{article}
\usepackage{mathblatt}
% gesetzt von werkzeuge/setzer.py; Lösungen nach bau/bauregeln.md „Lösungen“.
\newcommand{\kennung}{S2L}
\begin{document}
\pfheftstil
\fancyfoot[C]{{\footnotesize\color{mbgrau}\kennung\hfill\thepage}}
\pfheftkopf{Prozentuale Veränderung}{Klasse 7 \textperiodcentered{} Lösungen}

\begin{pfloesung}{}
\lz{1b)}{$54\text{ kg}$}{$10\,\% = 6\text{ kg}$; $60\text{ kg} - 6\text{ kg}$; auf $90\,\%$}{}
\lz{c)}{$60\text{ m}$}{$50\,\% = 20\text{ m}$; $40\text{ m} + 20\text{ m}$; auf $150\,\%$}{}
\end{pfloesung}
\begin{pfloesung}{}
\lz{2b)}{$0{,}7$}{um $30\,\%$ weniger heißt auf $70\,\%$}{}
\lz{c)}{$0{,}1$}{um $90\,\%$ weniger heißt auf $100\,\% - 90\,\% = 10\,\%$}{}
\lz{d)}{$0{,}7$}{auf $70\,\%$: der neue Wert ist $70\,\%$ (nicht $0{,}3$)}{}
\lz{e)}{$1{,}05$}{auf $105\,\%$; $5\,\% = 0{,}05$}{}
\lz{f)}{$2$}{auf $200\,\%$: das Doppelte}{}
\end{pfloesung}
\begin{pfloesung}{}
\lz{3b)}{$72$ kg}{$90 \cdot 0{,}8$}{}
\lz{c)}{$204\,€$}{$240 \cdot 0{,}85$}{}
\lz{d)}{$3{,}78\,€$}{$3{,}50 \cdot 1{,}08$}{}
\lz{e)}{$1\,236$}{$1\,200 \cdot 1{,}03$}{}
\lz{f)}{$51\,€$}{$85 \cdot 0{,}6$; auf $60\,\%$, nicht um}{}
\end{pfloesung}
\begin{pfloesung}{}
\lz{4b)}{$15\,\%$ weniger}{Unterschied $12$; $12 : 80 = 0{,}15$}{}
\lz{c)}{$20\,\%$ mehr}{Unterschied $10$; $10 : 50 = 0{,}2$}{}
\lz{d)}{$20\,\%$ weniger}{Unterschied $10$; $10 : 50 = 0{,}2$; wie a), aber alt ist $50\,€$}{}
\end{pfloesung}
\begin{pfloesung}{}
\lz{5b)}{$100 : 500$}{alt ist früher: $500$, auch wenn es hinten steht}{}
\lz{c)}{$1{,}50 : 7{,}50$}{alt ist vorher: $7{,}50\,€$}{}
\lz{d)}{$200 : 1\,150$}{Mai ist vor Juni: alt $= 1\,150$}{}
\lz{e)}{$2 : 8$}{alt $= 6 + 2 = 8\,€$}{}
\end{pfloesung}
\begin{pfloesung}{}
\lz{6b)}{$66{,}96\,€$}{neuer Wert: $72 \cdot 0{,}93$}{}
\lz{c)}{$16{,}2\,\%$}{Veränderung: $550 : 3\,400 = 0{,}1617\ldots$}{}
\lz{d)}{$8{,}8\,\%$}{Veränderung: $1{,}1 : 12{,}5$}{}
\lz{e)}{$2{,}80\,€$}{neuer Wert: $2{,}50 \cdot 1{,}12$}{}
\lz{f)}{$8{,}6\,\%$}{Veränderung: $0{,}30 : 3{,}49 = 0{,}0859\ldots$}{}
\end{pfloesung}
\begin{pfloesung}{}
\lz{7a)}{$8{,}6\,\%$}{$0{,}15 : 1{,}75 = 0{,}0857\ldots$; alt ist Montag}{}
\lz{b)}{$4{,}2\,\%$}{$0{,}08 : 1{,}90 = 0{,}0421\ldots$; alt ist Freitag}{}
\lz{c)}{nein}{$0{,}08 : 1{,}82 \approx 4{,}4\,\%$ mehr; $0{,}08 : 1{,}90 \approx 4{,}2\,\%$ weniger; der alte Preis ist verschieden}{}
\end{pfloesung}
\end{document}
